%% ==========================================================================
%% docs/model_math.tex
%% Formal mathematical analysis of the Tail Assignment MILP:
%%   (1) Basic model (Khaled et al. 2018, Sections 3–4)
%%   (2) Maintenance-extended model (Section 5)
%%   (3) Full A/B/C/D hierarchy extension (our contribution)
%%   (4) Identification and correction of Constraint (13)
%% ==========================================================================
\documentclass[11pt, a4paper]{article}
\usepackage{amsmath, amssymb, amsthm}
\usepackage{booktabs}
\usepackage{hyperref}
\usepackage{xcolor}
\usepackage{cleveref}

\theoremstyle{plain}
\newtheorem{proposition}{Proposition}
\newtheorem{lemma}{Lemma}
\newtheorem{theorem}{Theorem}
\theoremstyle{definition}
\newtheorem{definition}{Definition}
\newtheorem{example}{Example}
\theoremstyle{remark}
\newtheorem{remark}{Remark}

\newcommand{\F}{\mathcal{F}}        % flight set
\newcommand{\P}{\mathcal{P}}        % aircraft set
\newcommand{\A}{\mathcal{A}}        % airport set
\newcommand{\D}{\mathcal{D}}        % day set
\newcommand{\C}{\mathcal{C}}        % check type set {A,B,C,D}
\newcommand{\MA}{\mathcal{M}}       % maintenance airport set
\newcommand{\Fl}[2]{F^{\downarrow}_{#1}(#2)}   % arrivals at k before theta-Delta
\newcommand{\Fd}[2]{F^{\uparrow}_{#1}(#2)}     % departures from k before theta
\newcommand{\aI}[1]{a^{\downarrow}_{#1}}        % arrival airport of flight i
\newcommand{\aD}[1]{a^{\uparrow}_{#1}}          % departure airport of flight i
\newcommand{\tI}[1]{t^{\downarrow}_{#1}}        % arrival time of flight i
\newcommand{\tD}[1]{t^{\uparrow}_{#1}}          % departure time of flight i
\newcommand{\dt}{\Delta t}                       % minimum ground time

\title{Mathematical Analysis of the Compact Tail Assignment Model\\
  \large Derivations, Proofs, and Constraint~(13) Correction}
\author{Working Document}
\date{\today}

\begin{document}
\maketitle
\tableofcontents
\clearpage

%%==========================================================================
\section{Notation Summary}
%%==========================================================================

\begin{table}[h]
\centering
\caption{Index sets and parameters.}
\begin{tabular}{ll}
\toprule
Symbol & Meaning \\
\midrule
$\F = \{1,\ldots,n_f\}$  & Set of flight legs \\
$\P = \{1,\ldots,n_p\}$  & Set of aircraft (tail numbers) \\
$\A = \{1,\ldots,n_k\}$  & Set of airports \\
$\D = \{1,\ldots,n_d\}$  & Set of planning days \\
$\C = \{A,B,C,D\}$       & Set of maintenance check types \\
$\MA \subseteq \A$        & Maintenance-capable airports \\
$\aD{i}$, $\aI{i}$       & Departure / arrival airport of flight $i$ \\
$\tD{i}$, $\tI{i}$       & Departure / arrival time of flight $i$ (in minutes) \\
$\dt$                     & Minimum ground turn time (minutes) \\
$c_{ij}$                  & Cost of assigning flight $i$ to aircraft $j$ \\
$a^0_j$                   & Initial position airport of aircraft $j$ \\
$T_{\max}^c$              & Max cumulative flight time before check $c \in \{A,B\}$ (min) \\
$d_{\max}^c$              & Max calendar days between consecutive checks $c \in \{C,D\}$ \\
$\delta^c$                & Duration of check $c$ (minutes) \\
$\text{mcap}_{m,d}$       & Maintenance capacity of airport $m$ on day $d$ (aircraft/day) \\
$\Phi^c_j$                & Pre-horizon accumulated time/days for aircraft $j$ toward check $c$ \\
$M_{dd'}$                 & Big-M for day pair $(d,d')$: max possible cumulative flight time \\
\bottomrule
\end{tabular}
\end{table}

\begin{table}[h]
\centering
\caption{Decision variables.}
\begin{tabular}{lll}
\toprule
Variable & Domain & Meaning \\
\midrule
$x_{ij}$      & $\{0,1\}$ & 1 iff flight $i$ is operated by aircraft $j$ \\
$z_{ijdc}$    & $\{0,1\}$ & 1 iff aircraft $j$ performs check $c$ after flight $i$ on day $d$ \\
$y_{jdc}$     & $\{0,1\}$ & 1 iff aircraft $j$ undergoes check $c$ on day $d$ \\
$\gamma_{jdc}$& $\{0,1\}$ & 1 iff aircraft $j$ has a check of level $\geq c$ on day $d$ \\
\bottomrule
\end{tabular}
\end{table}

\paragraph{Synchronization note (event-based update target).}
This document still presents the day-indexed maintenance extension for
historical continuity. The active improvement direction is the event-based
reformulation documented in \texttt{docs/update\_improve.tex}, where
maintenance is modeled as fixed-duration self-loop events with variable
$z_{ijk}$ (check $k$ after flight $i$ for aircraft $j$), and day-indexed
maintenance variables are removed.

%%==========================================================================
\section{Basic Model (Constraints C1--C4)}
%%==========================================================================

\subsection{Objective and Auxiliary Sets}

Define, for any airport $k \in \A$ and time threshold $\theta$:
\begin{align}
  \Fl{k}{\theta} &\;=\; \{i \in \F \mid \aI{i} = k,\; \tI{i} \leq \theta - \dt\} \\
  \Fd{k}{\theta} &\;=\; \{i \in \F \mid \aD{i} = k,\; \tD{i} < \theta\}
\end{align}

The \emph{net balance} of aircraft $j$ at airport $k$ up to time $\theta$ is:
\begin{equation}
  B(j, k, \theta) \;=\;
    \sum_{i \in \Fl{k}{\theta}} x_{ij}
    \;-\; \sum_{i \in \Fd{k}{\theta}} x_{ij}.
\end{equation}

\subsection{MILP Formulation}

\begin{align}
  \min_{x} \quad & \sum_{i \in \F} \sum_{j \in \P} c_{ij}\, x_{ij} \label{eq:obj} \\[6pt]
  \text{s.t.} \quad
  & \sum_{j \in \P} x_{ij} = 1, \quad \forall i \in \F \label{eq:c1} \\[4pt]
  & \sum_{i' \in \Fl{k}{\tD{i}}} x_{i'j}
    - \sum_{i' \in \Fd{k}{\tD{i}}} x_{i'j}
    \;\geq\; x_{ij}, \quad
    \forall j \in \P,\; k \in \A \setminus \{a^0_j\},\; i \in F^{\uparrow}_k
    \label{eq:c2} \\[4pt]
  & \sum_{i' \in \Fl{k}{\tD{i}}} x_{i'j}
    - \sum_{i' \in \Fd{k}{\tD{i}}} x_{i'j}
    \;\geq\; x_{ij} - 1, \quad
    \forall j \in \P,\; k = a^0_j,\; i \in F^{\uparrow}_k
    \label{eq:c3} \\[4pt]
  & x_{ij} \in \{0,1\}, \quad \forall i \in \F,\; j \in \P. \label{eq:bin}
\end{align}

Constraint~\eqref{eq:c1} ensures every flight is covered by exactly one aircraft.
Constraints~\eqref{eq:c2}--\eqref{eq:c3} together enforce equipment continuity
(C2), turn time (C4), and initial position (C3): a departure $i$ by aircraft $j$
at airport $k$ is feasible only if at least one prior arrival landed at $k$ (net
balance $\geq 1$ for non-home airports, $\geq 0$ for the home airport).

\begin{proposition}[Correctness; Khaled et al. 2018, Prop.~1]
  Any feasible solution to \eqref{eq:c1}--\eqref{eq:bin} is a feasible flight plan
  satisfying constraints C1--C4.
\end{proposition}
\begin{proof}[Proof sketch]
  Order aircraft $j$'s flights by departure time: $i_1, i_2, \ldots, i_p$.
  C3 (initial position): if $\aD{i_1} \neq a^0_j$, then $B(j, \aD{i_1}, \tD{i_1}) = 0$
  and $x_{i_1 j} = 1$ violates \eqref{eq:c3}—contradiction.
  C2 and C4: if the first pair $(i, i')$ violating continuity or turn-time is found,
  in both the non-home case ($k \neq a^0_j$) and home case ($k = a^0_j$) the
  respective net balance equals 0 or $-1$, violating \eqref{eq:c2} or \eqref{eq:c3}.
  Full proof by induction over the ordered flight sequence; see Khaled et al.
\end{proof}

\begin{proposition}[Compactness; Khaled et al. 2018, Prop.~2]
  The model \eqref{eq:obj}--\eqref{eq:bin} has $n_f \times n_p$ binary variables
  and $O(n_f \times n_p)$ constraints, where $n_f \gg n_p > n_k$.
\end{proposition}

%%==========================================================================
\section{Maintenance-Extended Model (Constraints C8--C15)}
%%==========================================================================

\subsection{Additional Notation}

Let $F(d, i)$ denote the set of flights $i'$ such that $\tI{i'} > \tI{i}$ and
$d_{i'} = d$ (later flights on the same day as $i$).
Let $F^{\downarrow}(m)$ be flights landing at maintenance airport $m \in \MA$.
Let $F_{d,d'} = \{i \in \F \mid d \leq d_i \leq d'\}$.

\subsection{Extended Objective}

\begin{equation}
  \min \sum_{i \in \F}\sum_{j \in \P} c_{ij}\, x_{ij}
       + \sum_{m \in \MA}\sum_{i \in F^{\downarrow}(m)}\sum_{j \in \P}\sum_{d \in \D}
         \text{mc}_{ijd}\, z_{ijdc}
  \label{eq:obj_ext}
\end{equation}

\subsection{Maintenance Constraints}

\paragraph{C8 — Maintenance blocks subsequent flights.}
\begin{equation}
  z_{ijdc} + x_{i'j} \leq 1,
  \quad \forall i \in \F,\; i' \in F(d,i),\; j \in \P,\; d \in \D,\; c \in \C
  \label{eq:c8}
\end{equation}

\paragraph{C9 — Maintenance requires assignment.}
\begin{equation}
  x_{ij} \geq z_{ijdc},
  \quad \forall i \in \F,\; j \in \P,\; d \in \D,\; c \in \C
  \label{eq:c9}
\end{equation}

\paragraph{C10 — Airport maintenance capacity.}
\begin{equation}
  \sum_{i \in F^{\downarrow}(m)} \sum_{j \in \P} z_{ijdc} \leq \text{mcap}_{m,d},
  \quad \forall d \in \D,\; m \in \MA,\; c \in \C
  \label{eq:c10}
\end{equation}

\paragraph{C11 — Aggregate $z$ to $y$.}
\begin{equation}
  \sum_{i \in F^{\downarrow}(m)} z_{ijdc} = y_{jdc},
  \quad \forall d \in \D,\; m \in \MA,\; j \in \P,\; c \in \C
  \label{eq:c11}
\end{equation}

\paragraph{C12 — Calendar-day spacing (C/D checks).}
For every window of $d_{\max}^c$ consecutive days:
\begin{equation}
  \sum_{r = d}^{d'} y_{jrc} \geq 1,
  \quad \forall j \in \P,\; c \in \{C,D\},\;
  d,d' \in \D \text{ s.t. } d' - d = d^c_{\max} - 1
  \label{eq:c12}
\end{equation}

%%==========================================================================
\section{Constraint (13): Error Identification and Correction}
%%==========================================================================

\subsection{Original Formulation (Khaled et al., eq.~13)}

The cumulative flight-hour constraint between two maintenance events for
hours-based checks ($c \in \{A, B\}$) reads in the original paper:

\begin{equation}
  \sum_{i \in F_{d,d'}} x_{ij}\, t_i
  \;\leq\;
  T^c_{\max}
  + M_{dd'}\bigl(2 - y_{jd c} - y_{jd'c}\bigr)
  + M_{dd'} \sum_{r = d+1}^{d'-1} y_{jrc},
  \label{eq:c13_original}
  \quad \forall d < d',\; 2 \leq d'-d \leq d^c_{\max}-1,\; j \in \P
\end{equation}

The intent is: if $y_{jdc}=1$ (check at start) AND $y_{jd'c}=1$ (check at end)
AND no check in between ($\sum_{r} y_{jrc} = 0$), then accumulated flight time
must not exceed $T^c_{\max}$.

\subsection{Identified Error}

\begin{lemma}[Constraint (13) is insufficient]
\label{lem:c13_error}
Constraint~\eqref{eq:c13_original} does \emph{not} enforce the flight-hour
limit between two successive maintenance events when one of the endpoint
indicators is zero.
\end{lemma}
\begin{proof}
Consider a scenario where aircraft $j$ has no check on day $d$ ($y_{jdc}=0$),
has a check on day $d'$ ($y_{jd'c}=1$), and no intermediate checks
($\sum_r y_{jrc}=0$). Substituting into \eqref{eq:c13_original}:
\[
  \sum_i x_{ij} t_i
  \;\leq\; T^c_{\max} + M_{dd'}(2 - 0 - 1) + 0
  \;=\; T^c_{\max} + M_{dd'}.
\]
Since $M_{dd'} \geq \sum_i t_i$ by construction, the constraint is trivially
satisfied regardless of how many flight hours accumulate. The flight-hour
limit is therefore \emph{not} enforced in this case, even though the interval
$[d, d']$ can represent the gap between two consecutive maintenance events
(the event at $d'$ being the next one after a previous event before $d$).
\end{proof}

\begin{remark}
The error arises because the big-M term is keyed on the \emph{product}
$y_{jdc} \cdot y_{jd'c}$ (via $2-y_{jdc}-y_{jd'c}$), which equals $M$ whenever
either endpoint is zero, rather than activating the constraint only when
both endpoints have scheduled checks.
\end{remark}

\subsection{Corrected Formulation}

\begin{equation}
  \sum_{i \in F_{d,d'}} x_{ij}\, t_i
  \;\leq\;
  T^c_{\max}
  + M_{dd'} \sum_{r=d+1}^{d'-1} y_{jrc}
  + M_{dd'}\, y_{jdc},
  \label{eq:c13_corrected_start}
\end{equation}
\begin{equation}
  \sum_{i \in F_{d,d'}} x_{ij}\, t_i
  \;\leq\;
  T^c_{\max}
  + M_{dd'} \sum_{r=d+1}^{d'-1} y_{jrc}
  + M_{dd'}\, y_{jd'c},
  \label{eq:c13_corrected_end}
\end{equation}
for all $j \in \P$, $c \in \{A,B\}$, $d,d' \in \D$ with $2 \leq d'-d \leq d^c_{\max}-1$.

\begin{proposition}[Correctness of the corrected formulation]
Constraints~\eqref{eq:c13_corrected_start}--\eqref{eq:c13_corrected_end}
correctly enforce $\sum_{i \in F_{d,d'}} x_{ij} t_i \leq T^c_{\max}$ whenever
$[d, d']$ represents an interval between two consecutive maintenance events
(no check in the interior, and no check at either endpoint that would signal
continuity to a neighboring interval).
\end{proposition}
\begin{proof}
Consider any pair $(d,d')$ and aircraft $j$. Three cases:
\begin{enumerate}
  \item $y_{jdc}=0$ and $y_{jd'c}=0$ (no check at either endpoint): in a
  feasible maintenance schedule, the interval $[d,d']$ is not a stand-alone
  maintenance interval. The big-M relaxation in both \eqref{eq:c13_corrected_start}
  and \eqref{eq:c13_corrected_end} makes both constraints vacuous — correct,
  since this interval is not the critical one.

  \item $y_{jdc}=1$ and $y_{jd'c}=1$ with $\sum_r y_{jrc}=0$ (checks at both
  endpoints, none in interior): both constraints reduce to
  $\sum x_{ij}t_i \leq T^c_{\max}$, enforcing the limit. $\checkmark$

  \item One endpoint has a check and the other does not (say $y_{jdc}=1$,
  $y_{jd'c}=0$): \eqref{eq:c13_corrected_start} is relaxed (big-M at $d$),
  but \eqref{eq:c13_corrected_end} becomes $\sum x_{ij}t_i \leq T^c_{\max}$
  (the $y_{jd'c}=0$ term vanishes), enforcing the limit.
  The symmetric case follows from \eqref{eq:c13_corrected_start}. $\checkmark$
\end{enumerate}
Together, the pair of constraints enforces the limit exactly for all intervals
that represent consecutive maintenance windows.
\end{proof}

\subsection{Big-M Computation}

For each day pair $(d, d')$, define:
\begin{equation}
  M_{dd'} = \max_{(i,i') \in F_d \times F_{d'}} \text{MFT}(i, i'),
\end{equation}
where $\text{MFT}(i, i')$ is the maximum total flight time along any feasible
sequence starting with flight $i$ (on day $d$) and ending with flight $i'$ (on day $d'$)
satisfying equipment continuity. This can be computed via a longest-path algorithm
on the acyclic connection graph (linear time in $|F_{d,d'}|$).

%%==========================================================================
\section{Full A/B/C/D Hierarchy Extension}
%%==========================================================================

\subsection{Hierarchy Variables}

The $\gamma_{jdc}$ (``mega-check'') variable indicates a check of level $\geq c$:
\begin{align}
  \gamma_{jdD} &= y_{jdD} \label{eq:hier_D} \\
  \gamma_{jdC} &= y_{jdC} + y_{jdD} \label{eq:hier_C} \\
  \gamma_{jdB} &= y_{jdB} + \gamma_{jdC} \label{eq:hier_B} \\
  \gamma_{jdA} &= y_{jdA} + \gamma_{jdB} \label{eq:hier_A}
\end{align}

Or equivalently: $\gamma_{jdc} = \sum_{c' \geq c} y_{jdc'}$.

\subsection{Reset Semantics}

When aircraft $j$ performs a D-check on day $d$, the reset rules are:
\begin{itemize}
  \item D-check resets $D, C, B, A$ accumulated days/hours.
  \item C-check resets $C, B, A$ (does not affect $D$ counter).
  \item B-check resets $B, A$.
  \item A-check resets $A$ only.
\end{itemize}

These semantics are enforced by using $\gamma_{jdc}$ (rather than $y_{jdc}$)
in constraints C12 and C13, so that any higher-level check counts as
satisfying the lower-level check requirement.

\subsection{Extended Constraint C13b: Pre-Horizon Accumulated Hours}

For the initial window from day 0 to day $d'$ ($d' < d^c_{\max}$), accounting
for hours already accumulated before the planning horizon begins:

\begin{equation}
  \sum_{i \in F_{0,d'}} x_{ij} t_i
  \;\leq\;
  (T^c_{\max} - \Phi^c_j)
  + M_{0,d'} \sum_{r=1}^{d'-1} \gamma_{jrc}
  + M_{0,d'}\, \gamma_{jd'c},
  \label{eq:c13b}
\end{equation}
for all $j \in \P$, $c \in \{A,B\}$, $1 \leq d' \leq d^c_{\max}-1$.

\subsection{Extended Constraint C14b: Multi-Day Check Duration}

If check $c$ has duration $\delta^c > 1440$ minutes (more than one day), it
occupies $\lceil \delta^c / 1440 \rceil$ consecutive days. No flight may be
assigned to aircraft $j$ on any of those days:

\begin{equation}
  y_{jdc} = 1 \implies
  \sum_{r=d}^{d + \lceil \delta^c/1440 \rceil - 1} y_{jr c} \geq
  \lceil \delta^c/1440 \rceil
  \label{eq:c14b}
\end{equation}

This is linearised as:
$\sum_{r=d}^{d + K_c - 1} y_{jrc} \geq K_c \cdot y_{jdc}$,
where $K_c = \lceil \delta^c / 1440 \rceil$.

%%==========================================================================
\section{Model Size Analysis}
%%==========================================================================

\begin{proposition}[Variable count]
The extended model has:
\[
  |\text{vars}| = n_f n_p + n_f n_p n_d |\C| + n_p n_d |\C| + n_p n_d |\C|
  \;=\; O(n_f n_p n_d |\C|).
\]
For $|\C|=4$ and typical $n_d = 30$, $n_f=1500$, $n_p=40$:
max $\approx 7.2 \times 10^6$ variables, but sparsity reduces active variables
by a factor of $>10\times$ (indexed only over valid flight-airport-day combinations).
\end{proposition}

\begin{proposition}[Constraint count]
The extended model has $O(n_f n_p n_d |\C|)$ constraints.
In practice (Table 10 of Khaled et al.): up to $\approx$130,000 constraints
for $(p=40, h=30)$.
\end{proposition}

%%==========================================================================
\section{LP Relaxation Strength}
%%==========================================================================

The LP relaxation of the basic model is remarkably tight:
from Table 8 of Khaled et al. (2018), the proportion of integral variables
in the LP optimal solution exceeds $85\%$ across all instances.
The integrality gap at the root node is below $0.5\%$ in all cases.

This can be partially explained by the following structural property:

\begin{proposition}[LP corner structure]
For any airport $k$ and aircraft $j$, the constraint matrix subblock
corresponding to constraints \eqref{eq:c2}--\eqref{eq:c3} is totally
unimodular (TU) if we ignore the off-diagonal terms coupling flights at
different airports. This near-TU structure produces very few fractional
LP solutions.
\end{proposition}

(Full proof requires analysis of the network flow substructure; see
Ahuja et al. \cite{ahuja1993network} for the TU result on flow matrices.)

\appendix

%%==========================================================================
\section{Counterexample for Original Constraint~(13)}
\label{app:c13_counterexample}
%%==========================================================================

\paragraph{Instance.}
Two flights, one aircraft ($j=1$), two days ($d=1$, $d'=2$), one airport.
$T^A_{\max} = 60$ minutes; $M_{12} = 200$.
\begin{itemize}
  \item Flight 1: duration 80 min, departs day 1.
  \item Flight 2: duration 70 min, departs day 2.
  \item No check on day 1 ($y_{1,1,A}=0$); check on day 2 ($y_{1,2,A}=1$).
\end{itemize}
Cumulative flight time: $80 + 70 = 150 > T^A_{\max} = 60$. Should be infeasible.

\paragraph{Original constraint \eqref{eq:c13_original}.}
\[
  150 \leq 60 + 200(2 - 0 - 1) + 0 = 60 + 200 = 260 \quad \checkmark \text{ (satisfied)}.
\]
Constraint is \emph{not} violated — incorrect.

\paragraph{Corrected constraints \eqref{eq:c13_corrected_end}.}
\[
  150 \leq 60 + 0 + 200 \times 0 = 60 \quad \times \text{ (violated — correct behaviour)}.
\]
The corrected constraint correctly identifies this as infeasible.

\bibliographystyle{plainnat}
\bibliography{../paper/refs}

\end{document}
